\scanpage{54}
\mspage{2}

La \deleted{complexe} \inserted{carte} \uline{duale} d'une carte (\deleted{\unsure{degré}} définie combinatoirement) a le même espace sous-jacent ``tri-angulé'' par les mêmes repères, mais les rôles des sommets de type $0$ et $\infty$ échangés: les sommets de $C$ correspondent aux faces de $\check C$ \deleted{\unsure{les arêtes de}} \inserted{et inversement} les faces de $C$ correspondent aux sommets de $\check C$, enfin les arêtes de $C$ et celles de $\check C$ se correspondent (avec conservation de la relation d'incidence). On notera que le repère commun $r$, \deleted{\unsure{muni}} quand il est \deleted{\unsure{muni}} regardé comme repère pour $\check C$ \deleted{\unsure{définie}} \inserted{en lieu de $C$}, reçoit des orientations \uline{opposées}, donc donne lieu à des orientations opposées sur les arêtes et pour les sens de rotation autour des sommets.

Une carte \inserted{$X$} est \uline{orientée} si la surface sous-jacente est orientée, ou ce qui revient au $=$, si dans l'ens $R(X)$ des repères de $X$, on a donné une subdivision en deux parties complémentaires
\[
\unsure{(5)|(6)}\qquad R(X)=R^+(X)\amalg R^-(X)
\]
telles que l'on ait
\[
(6\,\mathrm{bis})\qquad
\underset{\text{donc on a égalité}}{\unsure{\sigma_i|\tau_i}(R^+(X))\subset R^-(X)}
\quad\text{pour }i\in\{0,1,\infty\}
\]
ce qui implique
\[
\deleted{\unsure{(7)}}\quad
\unsure{\sigma_i|\tau_i}(R^+(X))=R^+(X),\quad
\unsure{\sigma_i|\tau_i}(R^-(X))=R^-(X)
\quad i\in\{0,1,\infty\}
\]
\[
(7\,\mathrm{bis})\qquad
\rho_i(R^+(X))=R^+(X),\quad\rho_i(R^-(X))=R^-(X)
\quad i\in\{0,1,\infty\}
\]
Alors \deleted{est} $R^+(X)$ est un $G_2^+$-ens à opérateurs, et
$X\unsure{\longrightarrow|{-}}R^+(X)$ donne une équiv. de catégories
\diagram{\begin{tikzcd}
(8)\quad\text{Cartes isot. or.}\arrow[r,"\approx"]&G_2^\circ\text{-ens }\inserted{\text{finis}}\text{ spéciaux}
\end{tikzcd}}
\margin{NB On utilise une notion de ``morphisme'' de cartes qui n'est pas tellement triviale (plus gén. que les iso)}

\mspage{53}
